\subsection[Studiengeb"uhren]{Studiengeb"uhren --\\ eine abschlie"sende
Betrachtung}
\emph{von Henning G"unther}

\begin{figure*}[b]
	\centering\includegraphics[width=0.6\textwidth]{bilder/comics/tigerente.png}
\end{figure*}

Wir schreiben das Wintersemester 2008/09.
Der Widerstand gegen Studiengeb"uhren liegt in Tr"ummern.
Nach den vernichtenden Niederlagen im voll\-st"an\-di\-gen Boykott der
Studiengeb"uhren im Sommersemster 2007, an dem nur 504 der "uber 14.000 Studenten teilnahmen und dem darauf folgenden, kaum noch sp"urbaren ,,\EUR{5}''-Boykott im Wintersemester 2007/08 sind die Studenten kaum noch zu Widerstand bereit. Im Sommersemster 2008 war das Werk vollbracht, jeder anf"angliche Widerstand in alle Winde zerstreut, die anf"anglich so breit erscheinende Front der Studiengeb"uhrengegner zerschlagen.

Was war geschehen?
Wie konnte sich die vormals so rebellische Studentenschaft, die fr"uher keine M"oglichkeit auslie"s, gegen das Unrecht zu protestieren, innerhalb von nur einem Jahr in einen in gedem"utigter Haltung die Geb"uhren entrichtenden Haufen Elend verwandeln?

Es hat den Anschein, dass die diabolisch geniale Saat der
Studiengeb"uhren-F"ursprecher, die Daumenschrauben der ,,Campus-Maut'' nicht
sofort und im vollen Umfang anzuziehen, auf ganzer Linie aufgegangen sei. Denn es traf zun"achst die, die sich am wenigsten wehren konnten: An Erstsemestern die, da noch nicht eingeschrieben, keinen Boykott wagen
konnten wurde zuerst erprobt, ob \EUR{500} ein Preis waren, f"ur den die Studenten zu k"ampfen bereit w"aren. Sie waren es nicht.
Zwar waren viele ,,im Prinzip'' dagegen, taten diese Meinung aber nur m"a"sig
auf den wenigen Demonstrationen kund.

Die meisten der Studenten scheinen sich inzwischen mit dem Fakt, mit j"ahrlich
\EUR{1000} weniger auskommen zu m"ussen, abgefunden zu haben. Kaum jemand gibt sich noch dem Wunschtraum hin, gr"o"sere Teile der Studenten f"ur irgendeine Form des organisierten Protest zu begeistern.
Es scheint fast als k"onnten die Studiengeb"uhrenschergen bald wieder Morgenluft wittern und in der Lage sein, dank mangelnden Widerstand, ihre k"uhnsten Tr"aume zu verwirklichen: \EUR{1000} Studiengeb"uhren pro Semester und mehr.

Was wird die Zukunft bringen?
Werden die Besiegten weiterhin wie die Gespenster einer l"angst vergangenen Zeit durch die Unig"ange huschen, von einer Vorlesung zur n"achsten hetzen, um sich durch ein schnelleres Studium vielleicht ein paar Euro Studiengeb"uhren zu sparen und gelernt haben, stets mit der Angst vor einer Erh"ohung der Geb"uhren zu leben?
Es bleibt zu hoffen dass den Advokaten des Bezahlstudiums dieser Triumph nicht gew"ahrt wird.
